\chapter{Background}\label{ch:background}

\section{Expected Goals in Football Analytics}

The quantitative analysis of goal-scoring probability in football has a long history. Early work by Reep and Benjamin~\cite{reep1968} attempted to model the relationship between field position and scoring outcomes, and later statistical studies formalised the idea that shot location is the dominant predictor of whether an attempt results in a goal. The modern Expected Goals (xG) metric grew out of this tradition: for a given shot, xG is the probability that it results in a goal, estimated from historical data as a function of distance to goal, angle to goal, shot type, and whether the chance was headed or from open play \cite{rathke2017}. A shot taken from six metres directly in front of goal might carry an xG of 0.70, whilst a speculative long-range attempt might carry 0.03.

xG has become a standard tool for football analysts and clubs because it provides a more stable measure of performance than raw goals scored \cite{statsbomb2018}. A team that consistently generates high-xG chances is likely to score more in the long run even if short-term results are noisy, and xG allows analysts to ask whether a result was deserved rather than simply lucky. Models have grown increasingly sophisticated, incorporating additional features such as the number of defenders between the shooter and the goal, body part used, and the nature of the build-up, with modern implementations trained on millions of shots from professional football \cite{lucey2015}.

However, all standard xG models share a fundamental limitation: they evaluate only shots. In a typical football match, shots account for fewer than 2\% of all on-ball events \cite{pappalardo2019}. The other 98\% (passes, carries, dribbles, and crosses) are ignored entirely by a shot-only xG framework. This means xG cannot tell us which pass created the high-quality chance, which carry advanced the ball into the danger zone, or which press forced the turnover that led to a goal. A metric that assigns a threat value to every possession event, not just shots, is therefore far more informative for understanding how football matches are actually decided.

\section{Quantifying Non-Shot Threat and Possession Value}\label{sec:bg_possession_value}

The most influential attempt to extend xG to non-shot events is Karun Singh's Expected Threat (xT) framework, introduced in 2019~\cite{singh2019}. Singh's formulation models the pitch as a grid of discrete zones and estimates, for each zone, the probability that a team in possession at that location will score within the next few actions. This is framed as a Markov chain: from any zone, the team can either shoot (with some shot probability and associated xG), move the ball to another zone (via a pass or carry), or lose possession. Iterating this recursively gives a grid of threat values that increase smoothly from deep defensive positions to the penalty area. The Markov-chain treatment of possession value has antecedents in earlier work on field-position value in football, but Singh's formulation popularised it as a practitioner tool and is the reference point for most subsequent non-shot threat models, including this one.

Formally, Singh's xT is defined by the following Markov-chain recursion over a discrete grid of $M$ zones. Let $V(s)$ denote the threat value of a team in possession at zone $s$. At each step, the team can either shoot (with probability $P_\text{shot}(s)$, yielding expected goals $G(s)$) or move the ball to another zone $s'$ (with probability $P_\text{move}(s \to s')$):
\[
    V(s) = P_\text{shot}(s)\cdot G(s) + P_\text{move}(s)\sum_{s'} T(s \to s')\cdot V(s')
\]
where $P_\text{move}(s) = 1 - P_\text{shot}(s)$ and $T(s \to s')$ is the empirical transition probability of moving from zone $s$ to zone $s'$ (estimated from historical data). This is a linear system in $V(s)$ that can be solved by value iteration. All quantities ($P_\text{shot}$, $G$, and $T$) are hand-computed from zone-level statistics. The result is a fixed lookup table.

Singh's xT grid is elegant and interpretable, but it has several limitations. The zone-based discretisation is coarse: two positions on opposite sides of a zone boundary receive identical threat values. The transition probabilities $T(s \to s')$ aggregate over all possible defensive configurations, so they cannot adjust when a particular zone is heavily defended. Most importantly, $V(s)$ depends only on the ball's location; the positions of the 21 other players are invisible to the model. This project takes Singh's core insight (that every possession event carries a measurable contribution to goal-scoring probability) and replaces the hand-crafted grid with a machine learning model trained end-to-end on raw event data. Rather than estimating $V(s)$ from zone-level transition counts, the model learns a continuous function of ball position, action type, and player context directly from labelled possession sequences, allowing it to capture finer spatial structure and player-position effects that the Markov-chain grid cannot represent.

A complementary, action-centric line of work values each action by the change in scoring and conceding probability it induces. The most cited example is Valuing Actions by Estimating Probabilities (VAEP), proposed by Decroos et al.~\cite{decroos2019}: rather than assigning a single threat value to a position, VAEP computes the change in goal probability and concession probability caused by each action, using a feature-engineered representation of the game state including recent event history. Decroos et al.\ report an ROC AUC of approximately 0.76 on an action-level goal-probability prediction task~\cite{decroos2019}; Chapter~\ref{ch:comparison} provides a controlled replication on the same StatsBomb test split used in this project (ROC AUC 0.7359), enabling a direct, within-dataset comparison with our three model versions. This credit-assignment philosophy was later given a reinforcement-learning treatment by Liu et al.~\cite{liu2020deep}, who learn an action-value ($Q$) function over play-by-play sequences with stacked LSTMs, framing player evaluation as estimating the long-run value of actions within a possession rather than the immediate next-goal probability. VAEP and the RL action-value models aim to attribute credit precisely to individual actions within a sequence, whereas xT aims to produce a spatial threat surface; our model is philosophically closer to the latter.

A third, spatially richer strand uses player-tracking data to model the value of \emph{space} and possession directly. Spearman's ``Beyond Expected Goals''~\cite{spearman2018} builds a physically motivated pitch-control field (the probability that each location is controlled by the attacking team given player positions, velocities, and plausible accelerations) and turns it into an off-ball spatial value model. Fernández and Bornn's ``Wide Open Spaces''~\cite{fernandez2018space} similarly quantifies how players create and occupy space, and Fernández, Bornn and Cervone's Expected Possession Value (EPV) framework~\cite{fernandez2019epv} unifies these ideas in a deep-learning model that estimates, at any instant, the expected value of a possession as a continuous function of the full spatial configuration. These models are the closest in spirit to what we attempt here: a continuous, player-aware value surface rather than a zone grid. The crucial practical difference is data. EPV and pitch-control models rely on continuous optical tracking (every player's position and velocity, many times per second), which is proprietary and unavailable in open form. We instead work from StatsBomb's openly released 360 freeze frames~\cite{statsbomb360}: a single static snapshot of visible player positions at the moment of each event, with no velocity and no off-screen players.

The present contribution learns a continuous threat surface (unlike Singh) and conditions on player positions (unlike Singh, Decroos, and the event-only RL models), but does so from sparse open freeze-frame snapshots rather than dense tracking (unlike Spearman and Fernández et al.). A secondary contribution, developed in Chapter~\ref{ch:v3}, is an explicit and honestly ablated test of whether token-level player attention adds value over an aggregate spatial encoding in this open-data regime, a question the tracking-data literature does not isolate, because it never lacks the dense spatial signal in the first place.

\section{Machine Learning in Sports Analytics}

The application of machine learning to football analytics has grown substantially in parallel with the availability of structured event data. Pappalardo et al.~\cite{pappalardo2019} demonstrated the research value of large-scale spatio-temporal event datasets by releasing a public dataset of match events from five European leagues, enabling reproducible work on tasks such as player performance rating, outcome prediction, and style classification. The availability of open data, including the StatsBomb open-data repository used in this project~\cite{statsbomb2018}, has since lowered the barrier to entry for this kind of research significantly.

Early machine learning approaches to football analytics typically relied on hand-crafted feature vectors: scalar summaries of ball position, action outcome, and game context fed into classical models such as logistic regression or gradient-boosted trees \cite{lucey2015}. These approaches work well for tasks where the relevant information can be expressed as a small number of interpretable numbers, but they struggle to capture the rich spatial structure of a football pitch, for example the difference between a pass into open space behind a high defensive line and an identical pass when the defence is compact and organised.

The trend in recent years has shifted towards architectures that learn spatial representations directly from data. Convolutional neural networks applied to pitch images, graph neural networks modelling player interactions, and Transformer-based models treating events as sequences have all been explored in the sports analytics literature \cite{yeung2023}. This project follows that trend: Version 1 uses a classical Random Forest on scalar features, Version 2 moves to a CNN that processes the pitch as a spatial image, and Version 3 introduces a Transformer to reason over individual player positions, each step corresponding to a richer and more expressive input representation.

\section{Convolutional Neural Networks for Spatial Data}

Convolutional Neural Networks (CNNs) were introduced by LeCun et al.~\cite{lecun1998} as an architecture designed to exploit the spatial structure of image data. Unlike fully-connected networks, which treat every input element independently, a CNN uses translation equivariance to apply small learnable filters across the input, allowing it to detect local spatial patterns regardless of where they appear. Stacking multiple convolutional layers allows the network to build up hierarchical representations, from simple edge detectors in early layers to complex semantic patterns in deeper layers.

CNNs have since found applications far beyond natural images. In satellite imagery, medical imaging, and physical simulation, any domain where the data has meaningful spatial structure and local patterns matter has benefited from convolutional architectures \cite{lecun1998}. The football pitch is an analogous domain: threat is spatially structured (the penalty area is dangerous, deep positions are not), and the relationship between nearby players is more important than interactions between distant ones. The network can learn that a dense cluster of defenders in the box suppresses threat, that an isolated attacker in space near goal amplifies it, and that the spatial pattern of play differs systematically by action type.

\section{Transformers and Attention Mechanisms}

The Transformer architecture, introduced by Vaswani et al.~\cite{vaswani2017}, replaced the recurrent structure of earlier sequence models with a mechanism called self-attention, in which every element of an input sequence can directly attend to every other element. For a sequence of $n$ tokens, self-attention computes a weighted sum of value vectors, where the weights are determined by the similarity between query and key vectors derived from the input. This allows the model to capture long-range dependencies without the vanishing gradient problems of recurrent networks, and to do so in parallel rather than sequentially.

Cross-attention extends this idea: queries come from one sequence whilst keys and values come from another. For Version 3, this means the Stage 1 threat estimate (logit) acts as the cross-attention query against the player context: given this baseline threat level, what adjustment do the player positions warrant?

The use of frozen pre-trained backbones combined with a small trainable head is well established in the transfer learning literature \cite{devlin2019}. Rather than training a Transformer on raw pitch data from scratch, Version 3 freezes the Version 2 CNN (which already encodes a strong positional threat estimate) and trains only the Stage 2 Transformer to produce a residual adjustment based on player positions. This keeps the training tractable and prevents the Transformer from simply re-learning what the CNN already knows. Treating each player as an individual token, rather than as a pixel in a density map, allows the model to reason about specific player relationships, for example which defender is closest to the ball or whether a goalkeeper is off their line, in a way that the CNN's convolutional filters cannot naturally represent.

\section{StatsBomb Open Data as a Research Resource}

StatsBomb is a football data company that has released a substantial portion of their event data as an open-access resource for academic and non-commercial use \cite{statsbomb2018}. The dataset covers over 3,400 matches across multiple competitions and leagues, with each match stored as a JSON file containing second-level timestamps, on-ball event types (passes, carries, shots, dribbles, and more), start and end coordinates, and outcome labels. It has become a standard dataset for football analytics research due to its breadth, consistent schema, and free availability via the StatsBombPy Python library \cite{statsbombpy}.

A more recent addition to the dataset is StatsBomb 360 data \cite{statsbomb360}, available for a subset of 425 matches. For each tracked event in a 360-compatible match, the dataset provides a freeze-frame snapshot of the positions of all visible players on the pitch at the moment of the event. This is qualitatively different from the standard event data: rather than knowing only where the ball is, the model can also see where every player is standing, enabling player-position-aware models. Versions 2 and 3 of this project rely on the 360 data to encode player context, whilst Version 1 uses only the standard event data available across the full 3,464-match dataset.
